\section{K-groups}
\label{section-K-groups}

\noindent
Let $R$ be a ring. We will introduce two abelian groups associated
to $R$. The first of the two is denoted $K'_0(R)$ and has the following
properties\footnote{The definition makes sense for any ring but is rarely
used unless $R$ is Noetherian.}:
\begin{enumerate}
\item For every finite $R$-module $M$ there is given an element $[M]$ in
$K'_0(R)$,
\item for every short exact sequence $0 \to M' \to M \to M'' \to 0$
of finite $R$-modules we have the relation $[M] = [M'] + [M'']$,
\item the group $K'_0(R)$ is generated by the elements $[M]$, and
\item all relations in $K'_0(R)$ among the generators $[M]$
are $\mathbf{Z}$-linear combinations
of the relations coming from exact sequences as above.
\end{enumerate}
The actual construction is a bit more annoying since one has to take care
that the collection of all finitely generated $R$-modules is a proper class.
However, this problem can be overcome by taking as set of generators
of the group $K'_0(R)$ the elements $[R^n/K]$ where $n$ ranges over
all integers and $K$ ranges over all submodules $K \subset R^n$.
The generators for the subgroup of relations imposed on these elements
will be the relations coming from short exact sequences whose terms
are of the form $R^n/K$. The element $[M]$ is defined by
choosing $n$ and $K$ such that $M \cong R^n/K$ and putting
$[M] = [R^n/K]$. Details left to the reader.

\begin{lemma}
\label{lemma-length-K0}
If $R$ is an Artinian local ring then the length function
defines a natural abelian group homomorphism
$\text{length}_R : K'_0(R) \to \mathbf{Z}$.
\end{lemma}

\begin{proof}
The length of any finite $R$-module is finite,
because it is the quotient of $R^n$ which has finite length by
Lemma \ref{lemma-artinian-finite-length}. And the length function
is additive, see Lemma \ref{lemma-length-additive}.
\end{proof}

\noindent
The second of the two is denoted $K_0(R)$ and has the following
properties:
\begin{enumerate}
\item For every finite projective $R$-module $M$ there
is given an element $[M]$ in $K_0(R)$,
\item for every short exact sequence $0 \to M' \to M \to M'' \to 0$
of finite projective $R$-modules we have the relation $[M] = [M'] + [M'']$,
\item the group $K_0(R)$ is generated by the elements $[M]$, and
\item all relations in $K_0(R)$ are $\mathbf{Z}$-linear combinations
of the relations coming from exact sequences as above.
\end{enumerate}
The construction of this group is done as above.

\medskip\noindent
We note that there is an obvious map $K_0(R) \to K'_0(R)$
which is not an isomorphism in general.

\begin{example}
\label{example-K0-field}
Note that if $R = k$ is a field then we clearly have
$K_0(k) = K'_0(k) \cong \mathbf{Z}$ with the isomorphism
given by the dimension function (which is also the length function).
\end{example}

\begin{example}
\label{example-K0-PID}
Let $R$ be a PID. We claim $K_0(R) = K'_0(R) = \mathbf{Z}$.
Namely, any finite projective $R$-module is finite free.
A finite free module has a well defined rank by Lemma \ref{lemma-rank}.
Given a short exact sequence of finite free modules
$$
0 \to M' \to M \to M'' \to 0
$$
we have $\text{rank}(M) = \text{rank}(M') + \text{rank}(M'')$
because we have $M \cong M' \oplus M'$ in this case (for example
we have a splitting by Lemma \ref{lemma-lift-map}).
We conclude $K_0(R) = \mathbf{Z}$.

\medskip\noindent
The structure theorem for modules of a PID says that
any finitely generated $R$-module is of the form
$M = R^{\oplus r} \oplus R/(d_1) \oplus \ldots \oplus R/(d_k)$.
Consider the short exact sequence
$$
0 \to (d_i) \to R \to R/(d_i) \to 0
$$
Since the ideal $(d_i)$ is isomorphic to $R$ as a module
(it is free with generator $d_i$), in $K'_0(R)$ we have
$[(d_i)] = [R]$.  Then $[R/(d_i)] = [(d_i)]-[R] = 0$.  From this it
follows that a torsion module has zero class in $K'_0(R)$.
Using the rank of the free part gives an identification $K'_0(R) = \mathbf{Z}$
and the canonical homomorphism from $K_0(R) \to K'_0(R)$ is an isomorphism.
\end{example}

\begin{example}
\label{example-K0-polynomial-ring}
Let $k$ be a field. Then $K_0(k[x]) = K'_0(k[x]) = \mathbf{Z}$.
This follows from Example \ref{example-K0-PID} as
$R = k[x]$ is a PID.
\end{example}

\begin{example}
\label{example-K0-node}
Let $k$ be a field. Let $R = \{f \in k[x] \mid f(0) = f(1)\}$,
compare Example \ref{example-affine-open-not-standard}.
In this case $K_0(R) \cong k^* \oplus \mathbf{Z}$, but
$K'_0(R) = \mathbf{Z}$.
\end{example}

\begin{lemma}
\label{lemma-K0-product}
Let $R = R_1 \times R_2$. Then $K_0(R) = K_0(R_1) \times K_0(R_2)$
and $K'_0(R) = K'_0(R_1) \times K'_0(R_2)$
\end{lemma}

\begin{proof}
Omitted.
\end{proof}

\begin{lemma}
\label{lemma-K0prime-Artinian}
Let $R$ be an Artinian local ring.
The map $\text{length}_R : K'_0(R) \to \mathbf{Z}$
of Lemma \ref{lemma-length-K0} is an isomorphism.
\end{lemma}

\begin{proof}
Omitted.
\end{proof}

\begin{lemma}
\label{lemma-K0-local}
Let $(R, \mathfrak m)$ be a local ring. Every finite projective $R$-module
is finite free. The map $\text{rank}_R : K_0(R) \to \mathbf{Z}$
defined by $[M] \to \text{rank}_R(M)$ is well defined
and an isomorphism.
\end{lemma}

\begin{proof}
Let $P$ be a finite projective $R$-module. Choose elements
$x_1, \ldots, x_n \in P$ which map to a basis of $P/\mathfrak m P$.
By Nakayama's Lemma \ref{lemma-NAK} these elements generate $P$.
The corresponding surjection $u : R^{\oplus n} \to P$ has a splitting
as $P$ is projective. Hence $R^{\oplus n} = P \oplus Q$ with $Q = \Ker(u)$.
It follows that $Q/\mathfrak m Q = 0$, hence $Q$ is zero by
Nakayama's lemma. In this way we see that every finite projective
$R$-module is finite free.
A finite free module has a well defined rank by Lemma \ref{lemma-rank}.
Given a short exact sequence of finite free $R$-modules
$$
0 \to M' \to M \to M'' \to 0
$$
we have $\text{rank}(M) = \text{rank}(M') + \text{rank}(M'')$
because we have $M \cong M' \oplus M'$ in this case (for example
we have a splitting by Lemma \ref{lemma-lift-map}).
We conclude $K_0(R) = \mathbf{Z}$.
\end{proof}

\begin{lemma}
\label{lemma-K0-and-K0prime-Artinian-local}
Let $R$ be a local Artinian ring. There is a commutative
diagram
$$
\xymatrix{
K_0(R) \ar[rr] \ar[d]_{\text{rank}_R} & &
K'_0(R) \ar[d]^{\text{length}_R} \\
\mathbf{Z} \ar[rr]^{\text{length}_R(R)} & &
\mathbf{Z}
}
$$
where the vertical maps are isomorphisms by
Lemmas \ref{lemma-K0prime-Artinian} and \ref{lemma-K0-local}.
\end{lemma}

\begin{proof}
Let $P$ be a finite projective $R$-module. We have to show that
$\text{length}_R(P) = \text{rank}_R(P) \text{length}_R(R)$.
By Lemma \ref{lemma-K0-local} the module $P$ is finite free. So
$P \cong R^{\oplus n}$ for some $n \geq 0$. Then $\text{rank}_R(P) = n$ and
$\text{length}_R(R^{\oplus n}) = n \text{length}_R(R)$
by additivity of lengths (Lemma \ref{lemma-length-additive}).
Thus the result holds.
\end{proof}

















